% !TeX spellcheck = it_IT

\documentclass[handout]{beamer}

\usepackage[utf8]{inputenc}
\usepackage[italian]{babel}
\usepackage{multicol}
\usetheme{metropolis}
\usepackage{forest}

\usepackage{listings}
\usepackage{color} %red, green, blue, yellow, cyan, magenta, black, white
\definecolor{mygreen}{RGB}{28,172,0} % color values Red, Green, Blue
\definecolor{mylilas}{RGB}{170,55,241}
\lstset{language=Matlab,%
	%basicstyle=\color{red},
	breaklines=true,%
	morekeywords={matlab2tikz},
	keywordstyle=\color{blue},%
	morekeywords=[2]{1}, keywordstyle=[2]{\color{black}},
	identifierstyle=\color{black},%
	stringstyle=\color{mylilas},
	commentstyle=\color{mygreen},%
	showstringspaces=false,%without this there will be a symbol in the places where there is a space
	numbers=left,%
	numberstyle={\tiny \color{black}},% size of the numbers
	numbersep=9pt, % this defines how far the numbers are from the text
	emph=[1]{for,end,break},emphstyle=[1]\color{red}, %some words to emphasise
	%emph=[2]{word1,word2}, emphstyle=[2]{style},
	tabsize=2,
}



\usepackage{tikz}
\usetikzlibrary{shapes,arrows,positioning}

\tikzstyle{startstop} = [ellipse, draw, fill=bg, text=structure, text width=3em, text badly centered, inner sep=0pt]
\tikzstyle{block} = [rectangle, draw, fill=bg, text=structure, text width=5em, text centered]
\tikzstyle{decision} = [diamond, aspect=2, draw, fill=bg, text=structure, text width=3em, text badly centered, inner sep=0pt]
\tikzstyle{input} = [rectangle, draw, fill=bg, text=structure, text width=4em, text centered, rounded rectangle, rounded rectangle east arc=none]
\tikzstyle{output} = [rectangle, draw, fill=bg, text=structure, text width=4em, text centered, rounded rectangle, rounded rectangle west arc=none]

\tikzstyle{line} = [draw, -latex']


\title{Informatica B - Esercitazione 9}
\subtitle{Matlab - Funzioni ricorsive}

\author{Stefano Cereda\\
		stefano1.cereda@mail.polimi.it
	}
\date{5/12/2017}
\institute[PoliMi]{\vspace{0.5cm}\centering Politecnico di Milano \\ \vspace{0.2cm}
	\includegraphics[width=0.2\linewidth]{./logopolimi}}
%\logo{\includegraphics[width=15mm]{./logopolimi}}

\setbeamercovered{invisible}

\makeindex

\begin{document}
\begin{frame}
	\maketitle
	AGGIUNGERE BISEZIONE
\end{frame}

\begin{frame}[fragile]{Esercizio 0 - Struct e funzioni}
	Si consideri la seguente struttura dati Temperatura:
	\begin{lstlisting}
Temperatura.valore = 20;
Temperatura.scala = 'celsius'
	\end{lstlisting}
	Dove il campo \emph{scala} indica se il campo \emph{valore} è da intendersi espresso in gradi celsius o fahrenheit.
	
	Scrivere una funzione che, ricevuta in ingresso una struct Temperatura restituisca una struct in cui il valore di temperatura è espresso secondo l'altra scala.
	
	Si ricorda che: \[fahrenheit = celsius * (9/5) + 32;\] 	\[celsius = (fahrenheit -32) * (5/9)\]
\end{frame}

\begin{frame}{Esercizio 1 - Notazione scientifica}
	Scrivere una funzione iterativa che dato un numero reale positivo lo traduca in notazione scientifica.
	
	$123.456 = 1.23456E2$
	
	Scrivere la stessa funzione in modo ricorsivo.
\end{frame}

\begin{frame}{Esercizio 2 - Triangolo di Tartaglia}
	Si consideri una matrice triangolare T rappresentante il triangolo di Tartaglia, i cui valori sono definiti nel seguente modo:
	\begin{itemize}
		\item T(r,1) = 1
		\item T(r,r) = 1
		\item T(r,c) = T(r-1, c-1) + T(r-1, c)
	\end{itemize}

	Si scriva una funzione ricorsiva che, usando la definizione sopra, permetta di calcolare il valore di un generico elemento della matrice T a partire dagli indici di riga e colonna.
\end{frame}

\begin{frame}[fragile]{Esercizio 3 - Funzione misteriosa A}
	Si dica cosa calcola la seguente funzione ricorsiva quando si passano i parametri 7 e 2 ed in generale.
	
\begin{lstlisting}
function [z]=mistero1(x,y)
	if x<y
		z=0;
	else
		if mod(x,y) == 0
			z=x+mistero1(x-y,y);
		else
			z=mistero1(x-1,y);
		end
	end
	\end{lstlisting}
	
	\pause
	La funzione calcola la somma di tutti i multipli di $y$ compresi fra $x$ ed $y$.
\end{frame}

\begin{frame}{Esercizio 4 - Torri di Hanoi}
	Il gioco delle torri di Hanoi consiste nello spostare la pila di dischetti su un nuovo piolo, rispettando il vincolo di non mettere un disco grande su un disco piccolo.
	
	\pause
	Una pila di due dischi può essere spostata da un piolo $a$ ad un piolo $b$ spostando il primo disco sul piolo $c$, il secondo disco su $b$ ed infine il primo su $b$.
	
	\pause
	Una torre di $K$ dischi può quindi essere spostata su un nuovo piolo usando il terzo piolo come supporto intermedio:
	\begin{enumerate}
		\item sposto la pila di $k-1$ dischi sul piolo d'appoggio (ricorsione)
		\item sposto il primo disco sul piolo di destinazione (caso base)
		\item sposto la pila di $k-1$ dischi sul piolo di destinazione (ricorsione)
	\end{enumerate}
\end{frame}

\begin{frame}[fragile]{Esercizio 5 - Funzione misteriosa B}
	Si dica cosa calcola la seguente funzione ricorsiva quando si passa il parametro 7, quando si passa il parametro 8 ed in generale.
	
\begin{lstlisting}
function [z] = mistero2(x)
	if x>= 1
		z=mod(x,2) + 10*mistero2(floor(x/2))
	else
		z=0;
	end
\end{lstlisting}

\pause
mistero2(x) calcola un valore numerico in base 10 la cui sequenza di cifre può essere
interpretata come rappresentazione del numero naturale x nel sistema binario.
\end{frame}

\begin{frame}[fragile, allowframebreaks]{Esercizio 6 - Funzione misteriosa C - TdE 17/02/2009}
	Si consideri il seguente programma:
\begin{lstlisting}[basicstyle=\small]
function[ris] = s(n)
	if n<1
		ris = -1;
	elseif n>=1 & n<=4 %*
		ris = n; %*
	else
		ris = s(n-2)*s(n-4);
	end

%script di chiamata
for x = 1:8
	fprintf('%d, ', s(x));	
end	
\end{lstlisting}
	
	Quali risultati vengono stampati a video?
	
	Si supponga ora di rimuovere le linee evidenziate da un asterisco:
	\begin{enumerate}
		\item La terminazione della ricorsione è ancora garantita per ogni valore
		dell’argomento n? In caso affermativo giustificare brevemente la risposta,
		in caso negativo riportare almeno un esempio di argomento in cui la
		ricorsione non termina.
		
		\item Con la modifica apportata quali risultati vengono stampati a video?
	\end{enumerate}
\end{frame}

\begin{frame}[allowframebreaks]{Esercizio 7 - Ricorsione}
	Si scriva una funzione in MATLAB che abbia come parametro di ingresso un vettore contenente un insieme di valori, e che stampi tutte le permutazioni di
	tali valori.
	
	Le permutazioni di [1 2 3] sono:
	[1 2 3], [1 3 2], [2 1 3], [2 3 1], [3 1 2], [3 2 1].
	
	\begin{center}
	\begin{forest}
	[123
		[1
			[2
				[3]
			]
			[3
				[2]
			]
		]		
		[2
			[1
				[3]
			]
			[3
				[1]
			]
		]
		[3
			[1
				[2]
			]
			[2
				[1]
			]
		]	
	]
	\end{forest}
\end{center}

\framebreak
Osservando lo schema possiamo notare che abbiamo bisogno di due
informazioni: i valori già considerati precedentemente nelle permutazioni (per non ripeterli) e i valori ancora disponibili (da aggiungere). Queste informazioni dovranno quindi essere i parametri d'ingresso della funzione ricorsiva. Tuttavia, la consegna dice di scrivere una funzione che riceva come parametro solamente il vettore di valori da permutare, \alert{dovremo quindi aggiungere una funzione per mascherare il cambio di parametri}

\framebreak
La funzione si struttura nel seguente modo:
\begin{itemize}
	\item \alert{Caso base} L'insieme dei valori da inserire è vuoto: stampo i valori già inseriti e termino.
	\item \alert{Ricorsione} L'insieme dei valori da inserire non è vuoto: scandisco tutti i possibili valori
	ancora da inserire ed eseguo una invocazione ricorsiva della funzione per
	ognuno di questi valori. In queste invocazioni ricorsive passo come primo
	parametro i valori già inseriti dal chiamante, più quello nuovo che sto
	scandendo. Come secondo parametro invece passo tutti i valori ancora da
	inserire, eccetto quello che sto scandendo.
\end{itemize}
\end{frame}

\end{document}
